% Generated draft: requires mathematical and source review.
\documentclass[11pt]{article}
\usepackage[margin=25mm]{geometry}
\usepackage{fontspec}
\setmainfont{Latin Modern Roman}
\usepackage{amsmath,amssymb,mathtools,unicode-math}
\setmathfont{Latin Modern Math}
\usepackage{enumitem,needspace,xurl,hyperref}
\setlength{\parindent}{0pt}
\setlength{\parskip}{5pt}
\setlength{\emergencystretch}{3em}
\newcommand{\N}{\mathbb{N}}
\newcommand{\Z}{\mathbb{Z}}
\newcommand{\Q}{\mathbb{Q}}
\newcommand{\R}{\mathbb{R}}
\newcommand{\C}{\mathbb{C}}
\newcommand{\F}{\mathbb{F}}
\newcommand{\A}{\mathbb{A}}
\newcommand{\PP}{\mathbb{P}}
\newcommand{\E}{\mathbb{E}}
\newcommand{\eps}{\varepsilon}
\newcommand{\dd}{\,\mathrm{d}}
\begin{document}
\section*{UnsolvedMath ID: 30005722}
\textbf{Source problem code:} catalogue entry
\subsection{Definitions and mathematical statement}
The supplied record does not determine a standalone open problem. Its heading asks, ``How fast does it equilibrate?'', but no dynamics or particular equilibrium measure is specified. The record's verification note states that this question introduces log-Sobolev inequalities and quantitative equilibration bounds in the report.

The state space in the excerpt is \(X=\R^\Lambda\), with the index set \(\Lambda\) unspecified. For a probability measure \(\nu\) on \(X\) and a positive measurable function \(G\) with finite positive integral, define
\[
\operatorname{Ent}_\nu(G)=\int_X G\log\!\left(\frac{G}{\int_XG\,\mathrm d\nu}\right)\,\mathrm d\nu.
\]
For probability measures \(\mu\ll\nu\), with density \(f=\mathrm d\mu/\mathrm d\nu\), define relative entropy by
\[
H(\mu\mid\nu)=\int_X f\log f\,\mathrm d\nu
=\operatorname{Ent}_\nu(f),
\]
with \(0\log0=0\) and the value \(+\infty\) permitted. Define the total variation norm of the difference of two probability measures by
\[
\|\mu-\nu\|_{\mathrm{TV}}=2\sup_{A\text{ measurable}}|\mu(A)-\nu(A)|.
\]
This convention makes the displayed Pinsker bound below consistent; the excerpt does not specify its convention.

The first displayed formula in the supplied extraction suggests the log-Sobolev normalization
\[
\operatorname{Ent}_\nu(G)\leq\frac{1}{2\gamma}I_\nu(G),\qquad \gamma>0,
\]
for every bounded smooth strictly positive function \(G:\R^\Lambda\to(0,\infty)\). The excerpt calls the largest admissible \(\gamma\) the log-Sobolev constant. However, its extracted definition has displaced factors and does not unambiguously preserve the coefficient. Moreover, neither \(I_\nu\) nor the meaning of smoothness on \(\R^\Lambda\) is specified. For finite \(\Lambda\), a conventional Euclidean interpretation would be
\[
I_\nu(G)=\int_{\R^\Lambda}\frac{|\nabla G|^2}{G}\,\mathrm d\nu,
\]
but this interpretation is conditional and is not recoverable as an explicit definition from the supplied text.

Writing \(\nu_t\) for the unspecified time-dependent law and \(\nu_\infty\) for its intended equilibrium, the excerpt asserts
\[
H(\nu_t\mid\nu_\infty)
\leq\frac{1}{2\gamma}I(\nu_t\mid\nu_\infty)
=\frac{1}{2\gamma}I_{\nu_\infty}\!\left(\frac{\mathrm d\nu_t}{\mathrm d\nu_\infty}\right)
\]
and
\[
\|\nu_t-\nu_\infty\|_{\mathrm{TV}}^2
\leq2H(\nu_t\mid\nu_\infty)
\leq2e^{-2\gamma t}H(\nu_0\mid\nu_\infty),
\qquad t\geq0.
\]
The first inequality in the second display is attributed to Pinsker. The excerpt describes equilibration as having rate \(\gamma\); its display has exponent \(2\gamma\) for entropy and squared total variation. Absolute continuity is needed for the displayed densities, and finite initial entropy is needed for a finite bound. Deriving time decay also requires a specified evolution and an appropriate entropy-dissipation relation with compatible normalization; these hypotheses are absent. Consequently, these displays preserve the asserted quantitative content but do not constitute a fully specified theorem or a recoverable research question. No model-specific rate or sharp-constant problem can be stated without additional source information.
\subsection{Short English statement}
The excerpt presents log-Sobolev inequalities as a route to exponential equilibration, with an asserted entropy bound proportional to \(e^{-2\gamma t}\). It does not identify the dynamics, equilibrium measure, or a precise unresolved target. The supplied verification note identifies the passage as introductory exposition, although the registry labels it open. The extracted normalization and missing analytic hypotheses require human review.
\subsection{Sources}
\begin{itemize}
\item Local UnsolvedMath registry, permanent ID \(30005722\), title: ``Equilibration from Log-Sobolev Inequalities''; exact source problem code: \texttt{catalogue entry}. The supplied record belongs to the Oberwolfach Reports Open Problems collection and gives proposed year 2023.
\item Supplied report citation: \emph{Arbeitsgemeinschaft: QFT and Stochastic PDEs} (2024), Oberwolfach Reports, pp.\ 3319--3370, DOI \texttt{10.4171/owr/2023/59}. Supplied source and literature link, labeled ``Oberwolfach Report 59/2023'': \url{https://doi.org/10.4171/owr/2023/59}. The citation's 2024 designation and the report's 2023 designation are retained as supplied.
\item The extracted passage names Henri Elad Altman and Chunqiu Song. Their precise role is not specified, and the registry supplies no proposer attribution.
\item Literature assessment checked at \texttt{2026-08-22}: the supplied curation labels the item open and says that log-Sobolev theory gives exponential equilibration under a stated inequality, while the extracted question specifies neither a model nor a sharp rate. This is supplied dated curation, not independently verified literature assessment. The background expressly requests independent verification of time-sensitive status and cited claims.
\item The supplied statement status is \texttt{unrecoverable}. Its verification note says that the official report uses the question to introduce log-Sobolev inequalities and then supplies quantitative exponential equilibration bounds, rather than posing a standalone unresolved problem. This note is supplied evidence; the report has not been independently consulted for this draft.
\end{itemize}
\end{document}
